\section*{Capítulo \ref{ch:g7:plate-tectonics} --- La tectónica de placas: el rompecabezas que se mueve}

\begin{solution}{exo:g7:plate-tectonics:1}
La \omterm{def:g7:plate-tectonics:continental-drift}{deriva continental} es la idea de que los \omterm{def:g3:continents-and-oceans:continent}{continentes} se desplazan
lentamente sobre la superficie de la Tierra, de modo que sus posiciones han
cambiado a lo largo de las eras. Las pistas de Wegener: el encaje de las
\omterm{def:g1:landscapes-around-us:water}{costas} (América del Sur y África); los mismos fósiles en \omterm{def:g3:continents-and-oceans:continent}{continentes} hoy
separados por \omterm{def:g2:water-on-earth:ocean}{océanos}; las mismas \omterm{def:g3:stones-pebbles-sand:rock}{rocas} y antiguas \omterm{def:g4:world-physical-map:range}{cordilleras} a los dos
lados de un \omterm{def:g2:water-on-earth:ocean}{océano}; huellas de \omterm{def:g5:weather-and-climate:climate}{climas} antiguos fuera de lugar (carbón en la
Antártida, depósitos glaciares en la África cálida y seca).
\end{solution}

\begin{solution}{exo:g7:plate-tectonics:2}
Divergente: las \omterm{def:g7:plate-tectonics:plate}{placas} se separan. Convergente: se acercan. Transformante:
se deslizan una junto a otra, lado a lado.
\end{solution}

\begin{solution}{exo:g7:plate-tectonics:3}
Por ejemplo, todo alrededor del \omterm{def:g2:water-on-earth:ocean}{océano} Pacífico (Japón, Indonesia, los
Andes), y la franja que va del Mediterráneo al Himalaya. Casi no hay
grandes \omterm{def:g4:volcanoes-earthquakes-first-look:earthquake}{terremotos} en el interior de África ni de Australia (tampoco en
América del Norte al este de las \omterm{def:g1:landscapes-around-us:landform}{Montañas} Rocosas, ni en Brasil).
\end{solution}

\begin{solution}{exo:g7:plate-tectonics:4}
Un animal de \omterm{def:g2:water-on-earth:salt-fresh}{agua dulce} no podía cruzar miles de kilómetros de \omterm{def:g2:water-on-earth:ocean}{océano}
salado. Encontrarlo en los dos lados significa que los dos lugares
estuvieron unidos por tierra, en un mismo \omterm{def:g3:continents-and-oceans:continent}{continente}, cuando vivía
\emph{Mesosaurus}.
\end{solution}

\begin{solution}{exo:g7:plate-tectonics:5}
El carbón se forma con las plantas de bosques cálidos, húmedos y
pantanosos: el \omterm{def:g5:weather-and-climate:climate}{clima} era cálido y húmedo. Hoy esos bosques no pueden crecer
cerca del \omterm{def:g3:continents-and-oceans:poles}{Polo Sur}, así que la Antártida tuvo que estar mucho más cerca del
\omterm{def:g3:continents-and-oceans:equator}{ecuador}: se ha desplazado.
\end{solution}

\begin{solution}{exo:g7:plate-tectonics:6}
Dorsal mesoatlántica: divergente. Andes: convergente (una \omterm{def:g7:plate-tectonics:plate}{placa} oceánica se
hunde bajo un \omterm{def:g3:continents-and-oceans:continent}{continente}). Falla de San Andrés: transformante. Himalaya:
convergente (chocan dos \omterm{def:g3:continents-and-oceans:continent}{continentes}). Rift de África Oriental: divergente
(un \omterm{def:g3:continents-and-oceans:continent}{continente} que se parte).
\end{solution}

\begin{solution}{exo:g7:plate-tectonics:7}
$2.5 \times 10 = 25$ kilómetros por millón de años; en 60 millones de años:
$25 \times 60 = \qty{1500}{km}$.
\end{solution}

\begin{solution}{exo:g7:plate-tectonics:8}
$\qty{7}{m} = \qty{700}{cm}$, y $700 \div 9 \approx 78$: unos
\qty{78}{cm} al año, unas 30 veces la media de \qty{2.5}{cm} al año. Las
\omterm{def:g7:plate-tectonics:plate}{placas} no se separan suavemente en la superficie: la tensión se acumula
durante años y se libera en sacudidas breves, en las que las grietas se
abren varios metros de golpe. A lo largo de periodos largos, los años
tranquilos y las sacudidas se compensan y dan una media de unos pocos
centímetros al año.
\end{solution}

\begin{solution}{exo:g7:plate-tectonics:9}
No podía decir qué fuerza mueve los \omterm{def:g3:continents-and-oceans:continent}{continentes}, y el movimiento que
imaginaba (\omterm{def:g3:continents-and-oceans:continent}{continentes} que se abren paso a través del fondo del \omterm{def:g2:water-on-earth:ocean}{océano}) era
imposible. En la \omterm{def:g7:plate-tectonics:plate-tectonics}{tectónica de placas}, los \omterm{def:g3:continents-and-oceans:continent}{continentes} no se abren paso a
través de nada: son transportados, junto con el fondo oceánico que los
rodea, por \omterm{def:g7:plate-tectonics:plate}{placas} que se mueven como piezas enteras de la envoltura
exterior de la Tierra.
\end{solution}

\begin{solution}{exo:g7:plate-tectonics:10}
Más alto en la dorsal, más profundo en la fosa. Se fabrica \omterm{def:g7:plate-tectonics:plate}{placa} nueva en
la dorsal, donde sube \omterm{def:g3:stones-pebbles-sand:rock}{roca} caliente del \omterm{def:g6:inside-the-earth:layers}{manto}; la \omterm{def:g7:plate-tectonics:plate}{placa} oceánica vuelve al
\omterm{def:g6:inside-the-earth:layers}{manto} en la fosa, donde se hunde bajo el \omterm{def:g3:continents-and-oceans:continent}{continente}.
\end{solution}

\begin{solution}{exo:g7:plate-tectonics:11}
$14.2 \div 2.2 \approx 6.45$: unas 6 veces más rápida. En
kilómetros por millón de años, las velocidades son 142 y 22. Tiempos para
\qty{1000}{km}: $1000 \div 142 \approx 7$ millones de años para la dorsal
del Pacífico oriental; $1000 \div 22 \approx 45$ millones de años para la
dorsal del Atlántico Norte.
\end{solution}

\begin{solution}{exo:g7:plate-tectonics:12}
(a) $6472 - 4578 = 1894$ \omterm{def:g4:volcanoes-earthquakes-first-look:earthquake}{terremotos}. (b) $4578 \div 7031 \approx 0.651$,
es decir, el 65\,\%. (c) La línea del \omterm{def:g2:drawing-to-map:map}{mapa} marca la fosa, donde la \omterm{def:g7:plate-tectonics:plate}{placa}
oceánica empieza a hundirse. La \omterm{def:g7:plate-tectonics:plate}{placa} no baja en vertical: se hunde
inclinada bajo la otra \omterm{def:g7:plate-tectonics:plate}{placa}, cada vez más hondo a medida que avanza por
debajo de ella. Los \omterm{def:g4:volcanoes-earthquakes-first-look:earthquake}{terremotos} ocurren a lo largo de toda la \omterm{def:g7:plate-tectonics:plate}{placa} que se
hunde, así que muchos de ellos están bajo la \omterm{def:g7:plate-tectonics:plate}{placa} que queda encima, en el
mismo lado de la fosa, a decenas o cientos de kilómetros de la línea. Un
\omterm{def:g7:plate-tectonics:plate-boundary}{límite de placas} es una zona, no una línea fina.
\end{solution}

\begin{solution}{pb:g7:plate-tectonics:1}
\textbf{1.} \emph{Mesosaurus}, \emph{Cynognathus} y \emph{Glossopteris}
(la franja de \emph{Lystrosaurus} une África, la India y la Antártida).

\textbf{2.} Vivía en \omterm{def:g2:water-on-earth:salt-fresh}{agua dulce}: no podía cruzar un \omterm{def:g2:water-on-earth:ocean}{océano} salado de miles
de kilómetros de anchura.

\textbf{3.} La joroba del nordeste de Brasil encaja en el hueco del golfo de
Guinea, en la \omterm{def:g1:landscapes-around-us:water}{costa} oeste de África.

\textbf{4.} Un \omterm{def:g7:plate-tectonics:plate-boundary}{límite divergente}, la dorsal mesoatlántica: la \omterm{def:g7:plate-tectonics:plate}{placa}
Sudamericana se desplaza hacia el oeste y la \omterm{def:g7:plate-tectonics:plate}{placa} Africana hacia el este,
alejándose una de otra.

\textbf{5.} Las \omterm{def:g1:landscapes-around-us:water}{costas} están a lados opuestos del \omterm{def:g6:latitude-longitude:longitude}{meridiano de origen}:
$43.2 + 14.4 = \ang{57.6}$.

\textbf{6.} $57.6 \times 102 = 5875.2$: unos \qty{5880}{km}.

\textbf{7.} Los \omterm{def:g6:latitude-longitude:longitude}{meridianos} se acercan entre sí a medida que se va del
\omterm{def:g3:continents-and-oceans:equator}{ecuador} hacia un \omterm{def:g3:continents-and-oceans:poles}{polo}. El \omterm{def:g6:latitude-longitude:latitude}{paralelo} \ang{23} sur es un círculo más pequeño
que el \omterm{def:g3:continents-and-oceans:equator}{ecuador}, así que cada uno de sus 360 grados es más corto.

\textbf{8.} Más o menos la mitad de la anchura: $5880 \div 2 = \qty{2940}{km}$.

\textbf{9.} $3.4 \times 10 = 34$ kilómetros por millón de años.

\textbf{10.} $5880 \div 34 \approx 173$: unos 170 millones de años.

\textbf{11.} Sí: 170 millones de años está entre 100 y 200 millones de
años, y es posterior al comienzo de la fragmentación de \omterm{def:g7:plate-tectonics:supercontinent}{Pangea}, hace entre
225 y 200 millones de años. Un cálculo sencillo en un \omterm{def:g2:drawing-to-map:map}{mapa} da el orden de
magnitud correcto.

\textbf{12.} Por ejemplo: la velocidad no ha sido siempre la misma; los
\omterm{def:g3:continents-and-oceans:continent}{continentes} se estiraron y se adelgazaron antes de partirse, así que las
\omterm{def:g1:landscapes-around-us:water}{costas} actuales no coinciden exactamente con la línea por la que se
separaron.
\end{solution}
